\section{Mythos and Logos}

\begin{quotex}
We want to substitute faith for law, mythos for logos... will for pure reason, the image for the concept, and home for exile
\end{quotex}


Therein lays \textbf{Alain de Benoist}'s creed, which he takes seriously and every reader of \textbf{New Right} literature also needs to take seriously. When Benoist claims to reject \textbf{Charles Maurras} in toto, he necessarily includes Maurras' commitment to reason, which Benoist regards as a fault. In his interview, he makes this perfectly clear:

\begin{quotex}
Maurras is certain that it suffices to stop trusting in one's sentiments and to appeal to ``reason'', to avoid falling into utopian thinking and to reach to an almost perfect system of order. He boasts of building his system by a suite of logically irrefutable deduction, while he is only making his demonstrations in the abstract. His error, fundamentally, is to believe it suffices to reason correctly to reach the truth. It is to think, as Emmanuel Beau de Lomenie noted: ``that reason is one, that everyone reasoning correctly can only lead to the sole conclusion, which is inevitably true, and that consequently all doctrine which, in place of trusting in the intuitions of one's passions or feelings, will be concerned to reason correctly on the relationship of man with society, will be able, as to the problem of the political and social structure to establish, to admit as sane and just only one sole solution, that of Charles Maurras.'' There is in that a pretention that can be judged as exorbitant. There is also a certain naïveté, since his adversaries hold close to the same thing, that they believe that they also reason correctly, yet lead to opposite conclusions.
\end{quotex}


So according to Benoist, reasoning is useless in determining socio-political ends, and he prefers to rely on the intuitions of his passions. So instead of logically presenting his views, he prefers to present a narrative, or something to believe in that conforms to one's deepest feelings. He is able to gather followers who likewise believe in that narrative; these are his people, the like-minded. What about differences of opinion? Since they cannot be resolved by reason, then diversity is the necessary position to hold. Different groups will embrace different narratives, corresponding to their particular intuitions. Hence, his commitment to multiculturalism. Any attempt to embrace a single narrative, suitable to all, or different, peoples is considered by him to be a totalitarianism, the result of a monotheism as he understands the term.

Now we at Gornahoor embrace a different anthropology, viz., that man is a \textbf{rational animal}. Hence, we have a prior commitment to reason and logic. Nevertheless, there is a certain sympathy with Benoist, not in his methods, but rather in his objections. The Traditional conception of reason, however, is not that of the Enlightenment. Rather it embraces an intuition, but an intuition of higher things, which is quite different from that of the feelings. Furthermore, since Beauty is one of the transcendentals, we also embrace an intuition of feelings of a specific kind, which are of a different order than the lower passions.

To the objection that reasoners reach opposed conclusions, we agree with Benoit that this is due to a difference in first principles or postulates. The Traditional view, however, is that these first postulates are known through intuition. Hence, they are known directly or they are not; that is the nature of intuition. But relying on the intuitions of one's passions can hardly be the superior alternative.

The frustrating thing about Benoist is his lack of respect for metaphysical, logical, and historical precision. We ask whether a given belief is true or false, whether it conforms to Tradition or not. This is irrelevant to Benoist since, for him, a belief has only an instrumental value. For example, his understanding of monotheism is quite defective. So he never asks whether monotheism is true or not. Rather, he is concerned with the effects of accepting it as true.

Over time, we can deal with the specific errors, but for now, we are content with some examples exhibiting Benoist's creedal points. The main issue is that we don't recognize these dichotomies. Rather than an either-or, there is a both-and, properly understood.


\paragraph{Faith and Law}


Benoist rejects law in favor of faith, which he intends in the sense of conviction or fidelity. Hence, one does not follow an external law, but rather one's convictions to a belief. We also value fidelity, loyalty and conviction, but we also recognize the importance of the great pagan lawgivers, such as \textbf{Manu}, \textbf{Solon}, \textbf{Numa Pompilius}.

\paragraph{Mythos and Logos}


Cosmic order is foundational to Tradition, whether called Dharma in the East, or Logos in the West ... this is pagan, Western, and Traditional. Yet, myth is also part of Tradition. Since the Logos is beyond discursive thinking, it can only be fully expressed in terms of symbols, legends, rites, and myths. Yet, as Guenon repeatedly points out, symbolism is not irrational, yet it is beyond reason.

\paragraph{Will and Reason}


The relationship of will to reason: in some conceptions, reason is felt as a limitation. However, unlike a physical law, there is no compulsion to be reasonable. In any case, Tradition accepts the unity of being; hence, there can be no difference between the transcendentals of truth, good, and beauty. The will is oriented towards the good and the reason is oriented toward the truth. These must be one and the same, a view described by the New Right as totalitarian.

What is the alternative? If each group, or worse, each individual, is a law unto himself, each asserting his own will, then conflicts are inevitable. Without reason, there is no way to adjudicate competing claims. In Benoist's polytheistic system, for example, man is at the mercy of the gods, beseeching their favors whenever possible. For example, Aeneas was hindered by Juno, yet supported by Venus. This view does nothing to eliminate conflict.

\paragraph{Image and concept}


Once again the concept must not be misunderstood in rationalistic terms, as does Benoist. Rather, it is the idea, which is known directly and intuitively. The image is a lower form of knowing, for those who are incapable of thinking in metaphysical terms.

This misunderstanding is especially evident in the New Right's obsession with monotheism. Classical monotheism, in the West, is based on \textbf{Plato}, \textbf{Aristotle}, \textbf{Augustine}, and \textbf{Aquinas}. Thus, it has deep roots in pagan thought, and is far from a Semitic thought form. Moreover, as Guenon points out, the western understanding is equivalent to the notion of \textbf{Ishwara} in the Vedantic schools, again emphasizing its pagan origins.

Hence, Benoist's description of monotheism has nothing in common with Western tradition, which can only be understood conceptually, that is, transcendent to any imaginings. Since Benoist is committed to the image, his view of monotheism is absurdly anthropomorphic and is, in truth, unrecognizable to anyone steeped in the Tradition of classical theism. This topic will require extended treatment.


\flrightit{Posted on 2012-03-13 by Cologero}

\begin{center}* * *\end{center}

\begin{footnotesize}
\begin{sffamily}

\texttt{Cologero on 2012-03-21 at 00:50 said:}

Everything is perspective. Even the claim to reason is a sham, just a mask for the will-to-power. Yet we all feel compelled to promote our own perspectives as though they are absolute and come from God. So which perspective has the most power? Choose and exclude.

\hfill

\end{sffamily}
\end{footnotesize}

